% =========================
% 基础文档设置
% =========================
\documentclass[UTF8,a4paper,12pt]{ctexart}

% 页面设置
\usepackage{geometry}
\geometry{
  left=3cm,
  right=2.5cm,
  top=2.8cm,
  bottom=2.5cm
}

% =========================
% 字体设置
% =========================
\usepackage{fontspec}

% 英文、数字字体
\setmainfont{Times New Roman}

% 中文正文字体：宋体
\setCJKmainfont{SimSun}

% 中文无衬线字体：黑体
\setCJKsansfont{SimHei}

% 代码字体
\setmonofont{Consolas}

% =========================
% 常用宏包
% =========================
\usepackage{graphicx}
\usepackage{amsmath}
\usepackage{amssymb}
\usepackage{booktabs}
\usepackage{float}
\usepackage{xcolor}
\usepackage{listings}
\usepackage{setspace}
\usepackage{titlesec}
\usepackage{fancyhdr}

% =========================
% 正文格式设置
% =========================

% 正文字号：小四
\AtBeginDocument{\zihao{-4}}

% 正文首行缩进 2 字符
\setlength{\parindent}{2em}

% 段间距
\setlength{\parskip}{0pt}

% 行距
\setstretch{1.3}

% =========================
% 标题编号设置
% =========================

% 一级标题编号：一、二、三
\renewcommand{\thesection}{\chinese{section}}

% 二级标题编号：1. 2. 3.
\renewcommand{\thesubsection}{\arabic{subsection}}

% 三级标题编号：（1）（2）（3）
\renewcommand{\thesubsubsection}{（\arabic{subsubsection}）}

% =========================
% 标题格式设置
% =========================

% 一级标题：黑体，小三，加粗，左对齐
% 效果：一、实验目的及内容
\titleformat{\section}
  {\heiti\zihao{-3}\raggedright}
  {\thesection、}
  {0pt}
  {}

% 二级标题：黑体，四号，加粗，左对齐
% 效果：1. 概述
\titleformat{\subsection}
  {\heiti\zihao{4}\raggedright}
  {\thesubsection.}
  {0.5em}
  {}

% 三级标题：黑体，小四，加粗，左对齐
% 效果：（1）数据预处理
\titleformat{\subsubsection}
  {\heiti\zihao{-4}\raggedright}
  {\thesubsubsection}
  {0.5em}
  {}

% 标题间距
\titlespacing*{\section}
  {0pt}{1.0em}{0.5em}

\titlespacing*{\subsection}
  {0pt}{0.7em}{0.3em}

\titlespacing*{\subsubsection}
  {0pt}{0.5em}{0.2em}

% =========================
% 页脚页码设置
% =========================
\pagestyle{fancy}
\fancyhf{}
\fancyfoot[C]{-\thepage-}
\renewcommand{\headrulewidth}{0pt}

% =========================
% 代码格式设置
% =========================
\lstdefinestyle{reportcode}{
  basicstyle=\ttfamily\zihao{5},
  numbers=left,
  numberstyle=\ttfamily\zihao{-5},
  frame=single,
  breaklines=true,
  columns=fullflexible,
  keepspaces=true,
  showstringspaces=false,
  tabsize=4,
  backgroundcolor=\color{gray!8},
  keywordstyle=\bfseries,
  commentstyle=\itshape,
  captionpos=b
}

\lstset{
  style=reportcode
}

% =========================
% 正文开始
% =========================
\begin{document}

\section{实验目的及内容}

这里填写实验目的和实验内容。

例如：

本次实验的主要目的是掌握计算机视觉中图像读取、显示、预处理和特征提取的基本方法，熟悉 OpenCV 的基本使用流程，并通过实验理解相关算法的实现过程。

\section{实验原理及技术路线}

\subsection{概述}

这里填写实验背景和总体说明。

\subsection{实验原理}

这里填写实验涉及的主要原理。

例如，对于灰度化处理，可以将彩色图像中的 RGB 三个通道按照一定权重转换为单通道灰度图像：

\[
Gray = 0.299R + 0.587G + 0.114B
\]

\subsection{技术路线}

这里填写实验流程。

例如：

读取图像 $\rightarrow$ 图像预处理 $\rightarrow$ 特征提取 $\rightarrow$ 结果显示 $\rightarrow$ 分析实验结果。

\section{所用仪器、材料}

本次实验所用的软件环境和硬件环境如下：

\begin{itemize}
  \item 操作系统：Windows 11 / Ubuntu 22.04
  \item 编程语言：Python
  \item 开发环境：VS Code / PyCharm / Jupyter Notebook
  \item 主要库：OpenCV、NumPy、Matplotlib
\end{itemize}

\section{实验方法、步骤}

这里填写实验步骤。注意不要只粘贴代码，应当配合必要的解释。

\subsection{图像读取与灰度化处理}

首先使用 OpenCV 读取输入图像，然后将彩色图像转换为灰度图像，最后使用 Matplotlib 对处理结果进行可视化显示。

示例代码如下：

\begin{lstlisting}[language=Python]
import cv2
import matplotlib.pyplot as plt

img = cv2.imread("test.jpg")
gray = cv2.cvtColor(img, cv2.COLOR_BGR2GRAY)

plt.imshow(gray, cmap="gray")
plt.title("Gray Image")
plt.axis("off")
plt.show()
\end{lstlisting}

上述代码首先使用 OpenCV 读取图像，然后通过 \texttt{cv2.cvtColor} 将彩色图像转换为灰度图像，最后使用 Matplotlib 显示处理结果。

\subsection{结果显示与保存}

在完成图像处理后，可以进一步将结果保存到本地，便于后续分析和报告展示。

\begin{lstlisting}[language=Python]
cv2.imwrite("gray_result.jpg", gray)
\end{lstlisting}

该代码将灰度化后的图像保存为 \texttt{gray\_result.jpg} 文件。

\section{实验过程原始记录}

这里可以放实验截图、测试数据、运行结果、表格或中间计算过程。

例如：

\begin{table}[H]
\centering
\caption{实验过程记录表}
\begin{tabular}{ccc}
\toprule
实验编号 & 输入图像 & 处理结果 \\
\midrule
1 & test1.jpg & 成功完成灰度化 \\
2 & test2.jpg & 成功完成边缘检测 \\
\bottomrule
\end{tabular}
\end{table}

如果需要插入实验截图，可以使用如下格式：

\begin{figure}[H]
\centering
% \includegraphics[width=0.75\textwidth]{result.png}
\caption{实验结果示意图}
\end{figure}

注意：实际使用时，需要取消上面 \texttt{\textbackslash includegraphics} 这一行前面的注释，并将 \texttt{result.png} 替换成自己的图片文件名。

\section{实验结果、分析和结论}

这里填写实验结果总结、结果分析、存在的问题、改进方向和实验收获。

例如：

通过本次实验，完成了图像读取、灰度化处理和结果显示等基本操作。实验结果表明，OpenCV 可以较方便地完成常见的图像处理任务。在实验过程中，通过对图像处理前后结果的对比，可以更加直观地理解图像预处理在计算机视觉任务中的作用。

同时，本次实验也暴露出一些问题。例如，在处理不同尺寸、不同光照条件下的图像时，简单的灰度化和边缘检测方法可能无法稳定获得理想结果。后续可以进一步尝试图像滤波、边缘检测、特征匹配、深度学习模型训练等更复杂的计算机视觉方法，以提升实验的完整性和应用价值。

综上，本次实验加深了对计算机视觉基本处理流程的理解，为后续开展图像分类、目标检测和语义分割等任务奠定了基础。

\end{document}